\section{非平凡正例：可交换故障--修复系统}
\label{sec:fault-repair}

对每个 $N\ge2$，令
\[
X_N=\{0,1\}^N,
\]
其中 $x_i=1$ 表示第 $i$ 个部件故障。取 $u_N=\Id_{X_N}$，故操作性微观空间 $B_N=X_N$ 有 $2^N$ 个状态。每一步均匀选择一个部件：若它故障，则以概率 $a\in(0,1]$ 修复；若它正常，则以概率 $b\in[0,1]$ 故障；否则保持不变。记相应 kernel 为 $K_{N,a,b}$。任务只判断系统是否完全无故障：
\[
g_N(x)=\one_{\{|x|=0\}},
\qquad |x|:=\sum_{i=1}^N x_i.
\]

\begin{theorem}[故障数商的精确闭合与最小性]
\label{thm:fault-repair}
令 $q_N(x)=|x|$，$Y_N=\{0,1,\ldots,N\}$。则对每个 $a>0$、$b\in[0,1]$，$q_N$ 任务保真并精确闭合，诱导出生--死亡 kernel 为
\begin{align}
L_{N,a,b}(k,k-1)&=\frac{ak}{N},\nonumber\\
L_{N,a,b}(k,k+1)&=\frac{b(N-k)}{N},\label{eq:fault-repair-L}\\
L_{N,a,b}(k,k)&=1-\frac{ak+b(N-k)}{N},\nonumber
\end{align}
边界外概率取零。更强地，任意任务保真且对 $K_{N,a,b}$ 精确闭合的确定性商都必须区分不同故障数，因而最小状态数恰为
\[
c^0_{K_{N,a,b}}(g_N)=N+1.
\]
同一个 $q_N$ 对整个参数族 $a\in[a_-,1]$、$b\in[0,1]$（$a_->0$）同时成立，其中每个参数可有自己的 $L_{N,a,b}$。
\end{theorem}

\begin{proof}
从故障数 $k$ 出发，下一步减一、加一或不变的概率只依赖 $k$，直接计数即得式~\eqref{eq:fault-repair-L}，所以
$K_{N,a,b}Q_{q_N}=Q_{q_N}L_{N,a,b}$。且 $g_N(x)=\one_{\{q_N(x)=0\}}$，故任务保真。

为证最小性，设另一个任务保真精确闭合商把故障数 $k<\ell$ 的两个状态合并。精确闭合与任务保真意味着从这两个初态出发，任意未来时刻的任务分布都相同。但从 $k$ 个故障出发，前 $k$ 步依次选中每个故障部件并全部修复的概率至少为
\[
\frac{k!}{N^k}a^k>0;
\]
在这条路径上从未选中正常部件，所以 $b$ 不影响该下界。由 $\ell>k$ 个故障出发，每步至多修复一个部件，$k$ 步内不可能清零。二者在时刻 $k$ 的任务分布不同，矛盾。因此所有不同故障数必须分开，至少需要 $N+1$ 个状态；$q_N$ 达到该下界。
\end{proof}

这个例子同时给出指数级的操作性压缩与严格最小性：
\[
2=|O_N|<N+1<2^N=|B_N|,
\qquad
\frac{N+1}{2^N}\to0.
\]
对 $q_N$ 的分区 $\pi_N$，有 $D_{\bar\cK_N}(\pi_N)=0$，从而
\[
\Gamma_N\le\frac{N+1}{2^N}\to0.
\]
这里不能写成 $\Gamma_N=0$：联合指标还包含正的相对状态数。该模型也说明“宏观”不等于直接任务输出；完全无故障与非完全无故障只有两个任务标签，但为了自治预测，必须保留全部 $N+1$ 个故障数状态。
